\documentclass{situationist-preprint}

\title{Discovery and Invention in Mathematics}
\msid{MTH-03-a}
\msfamily{MTH-03}
\msdiscipline{Philosophy of Mathematics}
\mstype{Principal manuscript}
\mskeywords{discovery; invention; mathematical realism; notation; axiomatization; convergence; Wittgenstein}

\begin{document}
\maketitle

\begin{abstract}
Is mathematics invented or discovered? We argue that mathematical truths are discovered, while notations and axiomatizations are invented. The distinction is between what is chosen and what is constrained. A notation or an axiom system is chosen: alternatives were available and served equally well, or would have. Whether a sentence follows from given axioms, and whether it is true of a given structure, is not chosen: it is settled independently of anyone's decision, and mathematicians routinely find that it is settled otherwise than they expected. We support the thesis with three marks of discovery, namely independent convergence, resistance to the investigator, and the persistence of results across changes of notation and axiomatization, and we answer the objections from the freedom of mathematical creation and from Wittgenstein's view that a proof fixes a new rule rather than revealing a fact.
\end{abstract}

\section{The question}

Mathematicians speak of discovery and of invention, often of the same thing. \citet[§8]{Cantor1883} held that the essence of mathematics lies in its freedom, and \citet{Dedekind1888} described the numbers as free creations of the human mind. \citet{Hardy1940} held, on the contrary, that mathematical reality lies outside us and that the theorems we prove are observations of it. \citet{Wittgenstein1956} argued that the mathematician is an inventor, not a discoverer: a proof does not reveal a pre-existing fact but establishes a new criterion for the application of a concept.

We argue for a divided answer, and we locate the division precisely. Mathematical truths are discovered. Notations, definitions and axiomatizations are invented. The divided answer is not a compromise: it rests on a distinction between choices, which are free, and the consequences of choices, which are not.

\section{Discovery and invention}

\begin{definition}\label{d:inv}
An item is \emph{invented} if its existence or its form depends on a choice among alternatives that would have served the purposes for which it was introduced. An item is \emph{discovered} if it is found to obtain independently of the choices of those who find it.
\end{definition}

\begin{definition}\label{d:marks}
Three \emph{marks of discovery} are: (i) \emph{convergence}, independent inquirers reaching the same result from different starting points; (ii) \emph{resistance}, the result's being other than its investigators expected or wished; (iii) \emph{invariance}, the result's surviving changes in the invented apparatus by which it was reached.
\end{definition}

The marks are not definitions of discovery, but evidence for it. Each is what one would expect if a result is constrained independently of choice, and none is what one would expect of a choice.

\section{The argument}

\begin{proposition}\label{p:inv}
Notations and axiomatizations are invented.
\end{proposition}

\begin{proof}
Leibniz's notation for the calculus and Newton's served the same purposes, and the choice between them was settled by convenience. Arithmetic has been axiomatized by Dedekind's and Peano's axioms, in first-order and second-order versions, and by Hume's principle together with Frege's definitions (MTH-04); set theory has rival axiomatizations differing in their primitives and in whether they admit classes. In each case the alternatives were available, and each served or would have served. By Definition~\ref{d:inv}, these are inventions.
\end{proof}

\begin{proposition}\label{p:disc}
Whether a sentence follows from an axiom system, and whether it is true in an intended structure, is discovered.
\end{proposition}

\begin{proof}
Once the axioms and the logic are fixed, what follows from them is not further chosen: a theorem is derivable or it is not, and the mathematician who seeks a proof may fail to find one, or find a proof of the negation. Results exhibit all three marks. The calculus was developed independently by Newton and Leibniz, and non-Euclidean geometry independently by Gauss, Bolyai and Lobachevsky (convergence). Hilbert's programme expected a finitary consistency proof for arithmetic, and Gödel showed that none can be formalized in arithmetic itself (resistance). The infinitude of the primes holds whether the numbers are presented in decimal or binary notation, as von Neumann ordinals or as Hume's principle abstracts (invariance). By Definition~\ref{d:inv}, these results are discovered.
\end{proof}

\begin{theorem}\label{t:main}
Mathematics is discovered: mathematical truths are found, while notations and axiomatizations are invented.
\end{theorem}

\begin{proof}
Mathematical truths are truths about what follows from axioms and about what holds of mathematical structures, and these are discovered by Proposition~\ref{p:disc}. The apparatus by which they are stated and reached is invented by Proposition~\ref{p:inv}. That the apparatus is invented does not make the truths invented, since by the mark of invariance the truths survive changes in the apparatus.
\end{proof}

The theorem presupposes that there is something for mathematical truths to be true of. MTH-01 argues that numbers are abstract objects whose existence is established by the indispensability argument, and MTH-14 that they are positions in structures. The structures are not invented: a structure exists, on the ante rem view of MTH-14, whether or not anyone has described it, and an axiomatization is a description we choose. MTH-02 explains how such objects can be known, through reliable methods without causal contact. The companion manuscript MTH-03-b discusses the role of notation in discovery.

\section{Objections}

\begin{objection}[Freedom]
Mathematicians are free to define any concepts and adopt any consistent axioms they like; the complex numbers, the quaternions and non-Euclidean geometries were created, not found \citep{Cantor1883,Dedekind1888}.
\end{objection}

\begin{reply}
The freedom is a freedom of choice among structures to study, which is invention in the sense of Definition~\ref{d:inv}. It does not extend to what holds of the structures once chosen. Hamilton was free to look for an algebra of triples; he was not free to find one with the properties he wanted, and had to settle for quaternions, in which multiplication is not commutative. That is resistance, a mark of discovery.
\end{reply}

\begin{objection}[Rules]
A proof does not reveal what was already the case. It persuades us to adopt a new rule, under which we count a certain result as correct; before the proof, there was no fact \citep{Wittgenstein1956}.
\end{objection}

\begin{reply}
If proofs fixed new rules, there would be no explaining convergence: independent mathematicians adopting rules freely would be expected to diverge, as conventions in different communities do. Nor would there be any explaining why a proof can show a conjecture false that its proponents expected and wished to be true. The rule view also cannot account for the derivability of a theorem being settled before anyone has proved it, which is presupposed whenever a mathematician seeks a proof and when an unproved conjecture is used as a hypothesis whose truth value is unknown.
\end{reply}

\begin{objection}[Axiom choice]
If axiomatizations are invented, and truth is truth in an axiom system, then truth depends on invention after all. Different choices of axioms for set theory decide the continuum hypothesis differently.
\end{objection}

\begin{reply}
Truth in a system depends on the system chosen, and what is true in a given system is discovered. Whether the continuum hypothesis is true, rather than true in a system, depends on whether there is an intended structure that settles it. MTH-06 argues that it is false, given new axioms that are independently motivated; on that view the choice of new axioms is constrained by evidence about the intended structure, and is a matter of discovery guided by invention.
\end{reply}

\section{Conclusion}

Mathematical truths are discovered. Notations, definitions and axiomatizations are invented: they are chosen among alternatives that would have served. What follows from a choice, and what holds of the structure it describes, is not chosen, and it shows the marks of discovery: independent inquirers converge on it, it resists their expectations, and it survives the replacement of the apparatus by which it was found.

\begin{thebibliography}{9}\small
\bibitem[Cantor(1883)]{Cantor1883} Cantor, G. (1883). \emph{Grundlagen einer allgemeinen Mannichfaltigkeitslehre}. Teubner.
\bibitem[Dedekind(1888)]{Dedekind1888} Dedekind, R. (1888). \emph{Was sind und was sollen die Zahlen?} Vieweg.
\bibitem[Hardy(1940)]{Hardy1940} Hardy, G.~H. (1940). \emph{A Mathematician's Apology}. Cambridge University Press.
\bibitem[Wittgenstein(1956)]{Wittgenstein1956} Wittgenstein, L. (1956). \emph{Remarks on the Foundations of Mathematics}. Ed. G.~H. von Wright, R.~Rhees and G.~E.~M. Anscombe. Blackwell.
\end{thebibliography}

\end{document}
